\section{Proofs for Section~\ref{sec:acquirer}}\label{app:proofs-acquirer}

Numerical and symbolic checks of every statement below are in \texttt{code/verify\_part2\_acquirer.py} (sections A to D, F and G; section E checks Section~\ref{sec:acquired-growth}).

\paragraph{Proof of Proposition~\ref{prop:acq-owner}.}
(a) Write the stranded and rollback term of \eqref{eq:tasks-npv} as $[\varpi+(1-\varpi)p_u]C^{\rm str}+(1-\varpi)p_uL_{\rm roll}=p_u(C^{\rm str}+L_{\rm roll})+\varpi[(1-p_u)C^{\rm str}-p_uL_{\rm roll}]$. Then $\mathrm{NPV}_p=-(F+I+C^J)-C^{\rm par}A_r(T_{\rm par})+\bar b_p\varpi W\mathcal S_r-e^{-rT_{\rm par}}\{p_u(C^{\rm str}+L_{\rm roll})+\varpi[(1-p_u)C^{\rm str}-p_uL_{\rm roll}]\}$, which is affine in $F$, in $\varpi W\mathcal S_r$ and in $\varpi$. Evaluating it at $(F^{\mathsf o},W^{\mathsf o},\varpi^{\mathsf o},\mathcal S^{\mathsf o}_r)$ and at $(F,W,\varpi_p,\mathcal S_r)$, with all other inputs equal, and subtracting gives \eqref{eq:acq-channels}.

(b) Treat $N=N_{\rm pl}\ge1$ as real. By Proposition~\ref{prop:tasks-security}(a), $C^{\rm risk}$ is nondecreasing and concave with $C^{\rm risk}(0)=0$, so $C^{\rm risk}(x)\ge xC^{\rm risk\prime}(x)$ and $x\mapsto C^{\rm risk}(x)/x$ is nonincreasing. In \eqref{eq:acq-Wo}, $[C^{\rm pool}+\mathbb 1^{\rm reb}C^{\rm spec}]/N$ falls with $N$, and $C^{\rm risk}(NC^{\rm inc})/N=C^{\rm inc}\,C^{\rm risk}(NC^{\rm inc})/(NC^{\rm inc})$ is nonincreasing in $N$; hence $W^{\mathsf o}$ is nondecreasing in $N$. Moreover $W^{\mathsf o}|_{\mathbb 1^{\rm reb}=1}-W^{\mathsf o}|_{\mathbb 1^{\rm reb}=0}=C^{\rm spec}(1-1/N)\ge0$. Without a rebuild, $W^{\mathsf o}=W^0-C^{\rm spec}-[C^{\rm pool}+C^{\rm risk}(NC^{\rm inc})]/N<W^0-C^{\rm spec}$ whenever $C^{\rm pool}+C^{\rm risk}(NC^{\rm inc})>0$. With a rebuild, let $f(N):=NW^{\mathsf o}=NW^0-C^{\rm pool}-C^{\rm spec}-C^{\rm risk}(NC^{\rm inc})$. It is convex, because $C^{\rm risk}$ is concave. At $N=1$ the business bears all costs, $C^{\rm run}=C^{\rm pool}+C^{\rm spec}+C^{\rm risk}(C^{\rm inc})$, so $f(1)=W^0-C^{\rm run}<0$. If $C^{\rm risk}(x)/x\to0$ (under the exponential residual $C^{\rm risk}(x)=C_s[1+\ln(x/C_s)]$ for $x>C_s$), then $f(N)/N\to W^0>0$ and $f(N)\to\infty$. A convex function that is negative at 1 and tends to $+\infty$ has a sublevel set $\{N\ge1:f(N)\le0\}$ that is an interval containing 1; its right end $N^{\min}_{\rm pl}$ is the unique root above one, and $f>0$ beyond it. Finally, $f(C^{\rm run}/W^0)=C^{\rm risk}(C^{\rm inc})-C^{\rm risk}(C^{\rm inc}C^{\rm run}/W^0)\le0$ since $C^{\rm run}/W^0>1$, with strict inequality if $C^{\rm risk}$ is strictly increasing; so $N^{\min}_{\rm pl}\ge C^{\rm run}/W^0$, strictly in that case.

(c) Since $u_F\le1$ and $i^{-\epsilon_L}\le1$, $F^{\mathsf o}\ge F^{\rm rd}$; every subtracted term in \eqref{eq:acq-Wo} is nonnegative, so $W^{\mathsf o}\le W^0$. By (a), the value is decreasing in $F^{\mathsf o}$ and, with coefficient $\bar b_p\varpi^{\mathsf o}\mathcal S^{\mathsf o}_r\ge0$, nondecreasing in $W^{\mathsf o}$, which gives the floor. Monotonicity: $F^{\mathsf o}$ falls with $u_F$; $W^{\mathsf o}$ rises with $N_{\rm pl}$ and $\mathbb 1^{\rm reb}$ by (b); $\mathcal S^{\mathsf o}_r$ rises with $T^{\mathsf o}_h$ and enters with coefficient $\bar b_p\varpi^{\mathsf o}W^{\mathsf o}$, nonnegative if $W^{\mathsf o}\ge0$; and $\partial\mathrm{NPV}^{\mathsf o}_p/\partial\varpi^{\mathsf o}=\bar b_pW^{\mathsf o}\mathcal S^{\mathsf o}_r-e^{-rT_{\rm par}}[(1-p_u)C^{\rm str}-p_uL_{\rm roll}]$. If a capability gives $\mathrm{NPV}^{\mathsf o}_p>0$, any capability that is weakly larger in every coordinate, within the region where these signs hold, gives a value at least as large; so the flipping set is an upper set.

(d) In the frictionless benchmark of Proposition~\ref{prop:reg-lbo}(b) the new owner's NPV equals the value at $r_w$ of the firm under the plan it executes, less its outlay and expected distress costs; the argument there uses no property of the technology that produces the plan. Hence, for an owner with capability $\mathcal U_{\mathsf o}$ executing $\Pi^{\mathsf o}$, $\mathfrak p^{\max}_{\mathsf o}P_0=V^{\rm prog}(\Pi^{\mathsf o};\mathcal U_{\mathsf o})-V^{\rm prog}(\Pi^{\rm pub};\mathcal U_i)-C_{\rm tx}-e^{-r_wT}p_bL_b$, and for the sponsor the same with $V^{\rm prog}(\Pi^{\rm pe};\mathcal U_i)$. Subtracting gives \eqref{eq:acq-pmax}. The business's value to itself is $P_0=V^{\rm sq}+V^{\rm prog}(\Pi^{\rm pub};\mathcal U_i)$, and the owner's, net of costs and distress, is $V^{\rm sq}+V^{\rm prog}(\Pi^{\mathsf o};\mathcal U_{\mathsf o})-C_{\rm tx}-e^{-r_wT}p_bL_b$; comparing the two gives the stated condition, and comparing the owner's with the sponsor's gives $\Delta V^{\rm cap}_{\mathsf o}>0$. If the business is sold to the highest valuation at the second-highest price, the winner's surplus is the difference between the two valuations, which is zero when the runner-up holds the same capability.

(e) If $u_F$, $N_{\rm pl}$ and $\varpi^{\mathsf o}$ are nondecreasing in $i$, and $F^{\mathsf o}$ is nonincreasing in $i$ because $\epsilon_L\ge0$, then by (c) each process's $\mathrm{NPV}^{\mathsf o}_p$ is nondecreasing in $i$ where the sign conditions hold; $V^{\rm prog}$ is the sum of these values over the program, so $\Delta V^{\rm cap}_{\mathsf o}$, and by (d) $\mathfrak p^{\max}_{\mathsf o}$, are nondecreasing. If $m$ transformations run at once, each with at least $\underline n^{\rm core}$ FTE, then $m\underline n^{\rm core}\le\sum_in^{\rm core}_i\le N^{\rm core}$. \qed

\paragraph{Proof of Proposition~\ref{prop:acq-replace}.}
(a) Both paths share every flow except those listed in Assumption~\ref{ass:acq-replace}. In place, the $n_k$ processes of wave $k$ save $W$ a year from $t_k$, worth $n_kWe^{-rt_k}/r$; replacing, all $N=\sum_kn_k$ save from $t^{\rm rpl}$, worth $NWe^{-rt^{\rm rpl}}/r$. The difference is $\mathcal G^{\rm spd}$. Items (i) to (vi) are defined as present values at date 0, except the change in the floor, a flow of $S\Delta a_{\min}$ a year with present value $S\Delta a_{\min}/r$. So $V^{\rm rpl}-V^{\rm ip}=\mathcal G^{\rm spd}-\Delta\mathcal C$.

(b) $\partial\mathcal G^{\rm spd}/\partial t^{\rm rpl}=-NWe^{-rt^{\rm rpl}}<0$, and $t^{\rm rpl}$ is nondecreasing in each lead $T^{\rm leg}_{\mathfrak c}$. If $n'_{\rm cap}>n_{\rm cap}$, the $j$-th process embeds at $T_{\rm emb}+(\lceil j/n'_{\rm cap}\rceil-1)T\le T_{\rm emb}+(\lceil j/n_{\rm cap}\rceil-1)T$, so the in-place value rises and $\mathcal G^{\rm spd}$ falls. The terms of $\Delta\mathcal C$ are linear in severance, increasing in $\delta^G$, $T^{\rm cu}$ and $X^{\mathsf o}$ (through $C^{\rm cu}$) and in $L^{\rm ctx}$, and decreasing in $\kappa_{\rm doc}$ (through $C^{\rm ctx}$).

(c) With $n_k=n_{\rm cap}$ and $t_k=T_{\rm emb}+(k-1)T$, $\sum_{k=1}^Ke^{-rt_k}=e^{-rT_{\rm emb}}\sum_{k=0}^{K-1}e^{-rkT}=e^{-rT_{\rm emb}}(1-e^{-rKT})/(1-e^{-rT})$. As $r\to0$, $(e^{-rt^{\rm rpl}}-e^{-rt_k})/r\to t_k-t^{\rm rpl}$, so $\mathcal G^{\rm spd}\to W\sum_kn_k(t_k-t^{\rm rpl})=NW(\bar t^{\rm ip}-t^{\rm rpl})$.

(d) A year's cost is $L_{\rm mw}$ if at least one weakness is unremediated at its year-end, and cutovers create weaknesses independently given their probability. Spreading one cutover right after each of $n$ closes gives $nP_{\rm mw}(1-T_{\rm cl};1)$ expected weakness-years; bunching at $t^{\rm rpl}$ gives one year with probability $1-[1-P_{\rm mw}(1-z^{\rm rpl};n)]^n$. At equal positions ($z^{\rm rpl}=T_{\rm cl}$) and without congestion the two per-cutover probabilities are equal, and by Bernoulli's inequality, $(1-P)^n\ge1-nP$ for $P\in[0,1]$, so bunching is never worse. In general, bunching costs more if and only if $[1-P_{\rm mw}(1-z^{\rm rpl};n)]^n<1-nP_{\rm mw}(1-T_{\rm cl};1)$, which for $nP_{\rm mw}(1-T_{\rm cl};1)<1$ is \eqref{eq:acq-mw}. Equation \eqref{eq:acct-pmw} rises with $p_{\rm def}$, $T_{\rm dc}$ and $T_{\rm rem}$ and falls with $H$ (Proposition~\ref{prop:acct-cutover}(b)), so congestion and a later date each raise the left side. With a cost per weakness, bunching costs $nP_{\rm mw}(1-z^{\rm rpl};n)$ weakness-equivalents against $nP_{\rm mw}(1-T_{\rm cl};1)$, equal at equal positions without congestion and larger with congestion.

(e) Before $t^{\rm rpl}$ no process is embedded on the replacement path, and at $t^{\rm rpl}$ all are, with the permanent change $\delta\theta^{\rm rpl}$; the wave capacities of Assumption~\ref{ass:program} do not apply because the legacy organization does not absorb the change, so Proposition~\ref{prop:attainable}(c) does not bind. If every increment weakly enlarges the feasible set, then $\mathcal F(\theta_{\mathcal Z})\subseteq\mathcal F(\theta_{\mathcal P})$ for every $\mathcal Z\subseteq\mathcal P$, so $\mathcal A_t\subseteq\mathcal F(\theta_{\mathcal P})$ for all $t$, which is the replacement set at $t\ge t^{\rm rpl}$ when $\delta\theta^{\rm rpl}=0$. If $\delta\theta^{\rm rpl}$ lowers $R^*_X$, the in-place score $R^{\rm att}_X(t)$ is a nondecreasing step function (Proposition~\ref{prop:attainable}(a)) that reaches $R^*_X(\theta_{\mathcal P})>R^*_X(\theta_{\mathcal P}+\delta\theta^{\rm rpl})$ by $T_{\rm emb}+(K-1)T$ (Proposition~\ref{prop:attainable}(b)); let $t^\times$ be the first date at which it exceeds $R^*_X(\theta_{\mathcal P}+\delta\theta^{\rm rpl})$. Replacement attains at least the in-place score on $[t^{\rm rpl},t^\times)$ and strictly less afterward. All processes run in a single wave on one platform, and Proposition~\ref{prop:ext-portfolio}(c) bounds the probability that none succeeds below by $1-p_P$. \qed

\paragraph{Proof of Proposition~\ref{prop:acq-pc}.}
Let $T=T^{\rm IRR}$, $f(\mathfrak i):=\int_0^Te^{-\mathfrak it}\mathrm{FCF}\,dt+e^{-\mathfrak iT}\mathrm{TV}$ with $\mathrm{TV}=\int_T^\infty e^{-r_{\rm tv}(t-T)}\mathrm{FCF}\,dt$, and $g(\mathfrak i):=\int_0^\infty e^{-\mathfrak it}\mathrm{FCF}\,dt$. With $\mathrm{FCF}\ge0$ and not identically zero, both are strictly decreasing, and $f(r_{\rm tv})=g(r_{\rm tv})=V^{\rm tv}$.

(a) $f(\mathfrak i)-g(\mathfrak i)=e^{-\mathfrak iT}\int_T^\infty[e^{-r_{\rm tv}(t-T)}-e^{-\mathfrak i(t-T)}]\mathrm{FCF}\,dt$, which is nonnegative for $\mathfrak i\ge r_{\rm tv}$ and nonpositive for $\mathfrak i\le r_{\rm tv}$. Since $P=f(\mathfrak i^{\rm m})=g(\mathfrak i^\infty)$: if $\mathfrak i^\infty\ge r_{\rm tv}$, $f(\mathfrak i^\infty)\ge g(\mathfrak i^\infty)=f(\mathfrak i^{\rm m})$, so $\mathfrak i^{\rm m}\ge\mathfrak i^\infty$; if $\mathfrak i^\infty<r_{\rm tv}$, the same argument with the inequalities reversed gives $\mathfrak i^{\rm m}\le\mathfrak i^\infty$, strictly when flows after $T$ are positive. The right side of \eqref{eq:acq-method} is the IRR equation of a $T$-year hold whose sale price is the value at $r_{\rm tv}$ of the remaining flows, the market rule of Assumption~\ref{ass:reg-regime}(ii) with the buyer's rate $r_{\rm tv}$.

(b) Because $f$ is decreasing, $\mathfrak i^{\rm m}\ge\bar{\mathfrak i}$ if and only if $P=f(\mathfrak i^{\rm m})\le f(\bar{\mathfrak i})=P^{\max}(\bar{\mathfrak i})$. For $t\in[0,T]$ and $\bar{\mathfrak i}\ge r_{\rm tv}$, $e^{-r_{\rm tv}t}e^{-(\bar{\mathfrak i}-r_{\rm tv})T}\le e^{-\bar{\mathfrak i}t}\le e^{-r_{\rm tv}t}$, and $e^{-\bar{\mathfrak i}T}=e^{-r_{\rm tv}T}e^{-(\bar{\mathfrak i}-r_{\rm tv})T}$. Multiplying by $\mathrm{FCF}\ge0$ and $\mathrm{TV}\ge0$ and integrating gives \eqref{eq:acq-pmaxbound}. If $(1+\mathfrak p)V^{\rm tv}_{\rm sq}\le e^{-(\bar{\mathfrak i}-r_{\rm tv})T}(V^{\rm tv}_{\rm sq}+V^{\rm tv}_{\rm prog})$, then $P\le P^{\max}$ by the lower bound, which rearranges to the sufficient condition; and $P\le P^{\max}\le V^{\rm tv}$ gives $(1+\mathfrak p)V^{\rm tv}_{\rm sq}\le V^{\rm tv}_{\rm sq}+V^{\rm tv}_{\rm prog}$, the necessary one.

(c) Let the sweep pay the debt service $D(t)\ge0$, with any balance left at $T$ paid then, so that its present value at $r_{\rm debt}$ is $\mathrm{ND}_0$. Under the policy $\mathrm{ND}_0\le\mathrm{ND}^{\rm rep}$, and the balance $e^{r_{\rm debt}t}[\mathrm{ND}_0-\int_0^te^{-r_{\rm debt}s}\mathrm{FCF}\,ds]$ reaches zero by $T^{\rm rep}$, so if $T^{\rm rep}\le T$ nothing is left at $T$ and the equity flows after date 0, $\mathrm{FCF}-D$ and $\mathrm{TV}$, are nonnegative. Write $\mathrm{DS}(\mathfrak i)$ for the present value at $\mathfrak i$ of the debt service: it is strictly decreasing, with $\mathrm{DS}(r_{\rm debt})=\mathrm{ND}_0$. The equity value at $\mathfrak i$ less the equity outlay is $h(\mathfrak i):=f(\mathfrak i)-\mathrm{DS}(\mathfrak i)-(P-\mathrm{ND}_0)$, strictly decreasing with the unique root $\mathfrak i^{\rm m}_{\rm L}$. At the unlevered IRR, $f(\mathfrak i^{\rm m}_{\rm U})=P$, so $h(\mathfrak i^{\rm m}_{\rm U})=\mathrm{ND}_0-\mathrm{DS}(\mathfrak i^{\rm m}_{\rm U})$, which is positive if and only if $\mathfrak i^{\rm m}_{\rm U}>r_{\rm debt}$; and $h(\mathfrak i^{\rm m}_{\rm U})>0$ if and only if $\mathfrak i^{\rm m}_{\rm L}>\mathfrak i^{\rm m}_{\rm U}$. This is (i). For (ii), $\mathrm{ND}^{\rm rep}\le\overline{\mathrm{LTV}}P$ gives $\mathrm{ND}_0=\mathrm{ND}^{\rm rep}$, the sweep absorbs all free cash flow until $T^{\rm rep}=T$, and equity receives only $\mathrm{TV}$ at $T$: $P-\mathrm{ND}^{\rm rep}=e^{-\mathfrak iT}\mathrm{TV}$.

(d) In Proposition~\ref{prop:reg-weight}(c), $X^{\rm x}(\tau)=(1-e^{-(r-G)\tau})X^G+e^{-(r-G)\tau}X_T$ tends to $X^G=\hat m/(r-G)$ as $\tau\to\infty$; with $r=\bar{\mathfrak i}_{\rm U}$ and the dollar of revenue growing at $G_i$, $X^G=\hat m/(\bar{\mathfrak i}_{\rm U}-G_i)$. Dividing by the same expression at $r_w$ gives the ratio. \qed

\paragraph{Proof of Proposition~\ref{prop:acq-retention}.}
(a) By Proposition~\ref{prop:tasks-buildbuy}(c), a customer leaves rather than pay more than its switching cost per contract term, so it stays if and only if $\upsilon_{\rm rp}BT_c\le C^{\rm sw}$. Revenue retained is the new bill $(1+\upsilon_{\rm rp})B$ summed over customers who stay, $(1+\upsilon_{\rm rp})\mathsf S_{\rm sw}(\upsilon_{\rm rp})$ per dollar of previous revenue. With a single ratio $\bar x$, $\mathsf S_{\rm sw}(\upsilon)=1$ for $\upsilon\le\bar x$ and 0 above. The first-order condition of $\max_\upsilon(1+\upsilon)\mathsf S_{\rm sw}(\upsilon)$ is $\mathsf S_{\rm sw}+(1+\upsilon)\mathsf S'_{\rm sw}=0$; with $\mathsf S_{\rm sw}(x)=e^{-x/\bar x}$ it gives $\upsilon=\bar x-1$, a maximum since the objective is log-concave.
(b) Revenue this period from last period's customers is the sum over those who stay of their bills, which by definition of the revenue-weighted averages is $(1-\bar c)(1+\bar\upsilon)(1+\bar e)$ times last period's revenue. The penalty of Assumption~\ref{ass:ext-theta}(v) is a change in $G_i$, and by Proposition~\ref{prop:acquired-score}(d) $G_i=\ln\mathrm{NRR}$ for a business whose revenue is all subscription revenue, so the change is $\ln\mathrm{NRR}^{\rm pre}-\ln\mathrm{NRR}^{\rm post}$; a level of NRR at one date gives $G_i$ at that date only.
(c) Revenue at $t$ from the cohort of age $a$ is proportional to $e^{g^{\rm in}(t-a)}e^{-\nu^{\rm rv}a}$, so $\mathrm{Sh}(\tau)=\int_\tau^\infty e^{-(\nu^{\rm rv}+g^{\rm in})a}da/\int_0^\infty e^{-(\nu^{\rm rv}+g^{\rm in})a}da=e^{-(\nu^{\rm rv}+g^{\rm in})\tau}$ and the mean is $1/(\nu^{\rm rv}+g^{\rm in})$. A constant rate gives $\mathrm{Sh}(2\tau)=\mathrm{Sh}(\tau)^2$. If the rate $h$ varies across customers, $\mathrm{Sh}(\tau)=\mathbb E e^{-h\tau}$ and, by the Cauchy--Schwarz inequality, $\mathrm{Sh}(2\tau)=\mathbb E[(e^{-h\tau})^2]\ge(\mathbb Ee^{-h\tau})^2$, strictly unless $h$ is constant; the population hazard then falls with tenure as faster leavers exit.
(d) With $k=0$, $\hat X=X/(1-\tau)$ in \eqref{eq:ext-wedges}, so $\mathcal C$ is proportional to $X$; the ratio of the weights is that of Proposition~\ref{prop:acq-pc}(d). \qed

\paragraph{Proof of Proposition~\ref{prop:acq-limits}.}
(a) Per dollar of revenue, acquisitions use $k^{\rm acq}(g-g^{\rm org})$ and interest $r_{\rm debt}\mathrm{nd}$; the sources are $m$ and the new debt $d(\mathrm{nd}\,S)/dt/S=\mathrm{nd}\,g$ that keeps net debt at $\mathrm{nd}$ per dollar. Uses do not exceed sources if and only if $(k^{\rm acq}-\mathrm{nd})g\le m-r_{\rm debt}\mathrm{nd}+k^{\rm acq}g^{\rm org}$, which is $g\le g^{\max}$ since $k^{\rm acq}>\mathrm{nd}$; with $\mathrm{nd}=0$ it is $g\le g^{\rm org}+m/k^{\rm acq}$. Differentiating, $\partial g^{\max}/\partial k^{\rm acq}=-(m-r_{\rm debt}\mathrm{nd}+\mathrm{nd}\,g^{\rm org})/(k^{\rm acq}-\mathrm{nd})^2$. Since revenue grows at $g^{\rm org}+g^{\rm acq}$, holding $g$ requires $g^{\rm acq}=g-g^{\rm org}=g+d^{\rm org}$, which costs $k^{\rm acq}g^{\rm acq}$ of revenue. Without new customers $g^{\rm org}=G_i$ by Proposition~\ref{prop:acquired-score}(d), which is $\ln\mathrm{NRR}$ when all revenue is subscription revenue; at $g=0$, $g^{\rm acq}=-G_i$.
(b) Acquired revenue of $(g-g^{\rm org})S$ a year from at most $n^{\rm dl}$ deals needs a mean target of at least $(g-g^{\rm org})S/n^{\rm dl}$. Since $\mathsf h$ is decreasing for $\epsilon_{\rm sz}>0$, $\mathsf h(S^{\rm tg})\ge1$ if and only if $S^{\rm tg}\le S^{\rm tg}_{\max}$, so the hurdle admits at most $n^{\rm dl}S^{\rm tg}_{\max}$ of acquired revenue a year: $g-g^{\rm org}\le n^{\rm dl}S^{\rm tg}_{\max}/S$, which holds at the target $g$ if and only if $S\le S^{\max}$. If $\epsilon_{\rm sz}=0$ and $\mathsf h_0\ge1$, every target clears the hurdle.
(c) The $n$-th deal adds $V_n-P_n-p L^{\rm ch}$ in expectation, with $p=p^{\rm ch}(n\Delta s)$ under series review and $p^{\rm ch}(\Delta s)$ deal by deal. Since $V_n-P_n$ is nonincreasing and $p^{\rm ch}(n\Delta s)$ nondecreasing in $n$, the deals that pass form an initial segment, and $p^{\rm ch}(n\Delta s)\ge p^{\rm ch}(\Delta s)$ for $n\ge1$ makes the series segment no longer. With $p^{\rm ch}$ convex the increments $p^{\rm ch}((n+1)\Delta s)-p^{\rm ch}(n\Delta s)$ are nondecreasing. A review that delays a deal whose value and price both fall due at closing scales its net value by $e^{-rT^{\rm fdi}}$ with probability $p^{\rm fdi}$ and by one otherwise.
(d) A customer whose ratio is $h^{\rm cp}x$ stays after $\upsilon$ if and only if $h^{\rm cp}x\ge\upsilon$, so the retained share is $\mathsf S_{\rm sw}(\upsilon/h^{\rm cp})$, at most $\mathsf S_{\rm sw}(\upsilon)$ because $\mathsf S_{\rm sw}$ is nonincreasing and $h^{\rm cp}\le1$. With exponential ratios the scaled ratios are exponential with mean $h^{\rm cp}\bar x$, so by Proposition~\ref{prop:acq-retention}(a) the unconstrained optimum is $h^{\rm cp}\bar x-1$ when positive; the objective is log-concave, hence increasing below the optimum, and the cap binds when $\upsilon^{\rm cp}$ is lower. \qed
